\section{Implementasi Spesifikasi Wajib}\label{sec:Implementasi
Spesifikasi Wajib} % (fold)
\subsection{Algoritma Utama}\label{sub:Algoritma Utama} % (fold)

Algoritma Utama pada program ini mengimplementasikan strategi \textit{Divide and Conquer} melalui struktur data \textit{Octree} untuk melakukan proses \textit{voxelization}. Algoritma ini secara rekursif membagi ruang tiga dimensi menjadi delapan oktan hingga mencapai kedalaman maksimum yang ditentukan oleh pengguna. Untuk mendukung jalannya algoritma ini, digunakan beberapa struktur data utama, yaitu:

\begin{itemize}
    \item \texttt{OctreeNode} digunakan sebagai representasi setiap simpul dalam pohon \textit{octree}. Struktur ini menyimpan batas ruang (\textit{boundary}), referensi ke delapan anak simpul, serta status apakah simpul tersebut merupakan daun (\textit{leaf}) atau mengandung objek (\textit{solid}).
\begin{lstlisting}[language=Go, basicstyle=\ttfamily\small, breaklines=true]
type OctreeNode struct {
    Boundary BoundingBox
    Children [8]*OctreeNode
    IsLeaf   bool
    IsSolid  bool
}
\end{lstlisting}

    \item \texttt{BoundingBox} dan \texttt{Vector3} digunakan untuk mendefinisikan batas ruang dan koordinat dalam ruang 3D. \texttt{Vector3} menyimpan komponen $x, y, z$ dalam tipe \textit{float64}, sedangkan \texttt{BoundingBox} menyimpan titik minimum dan maksimum dari suatu kubus.
\begin{lstlisting}[language=Go, basicstyle=\ttfamily\small, breaklines=true]
type Vector3 struct {
    X, Y, Z float64
}

type BoundingBox struct {
    Min Vector3
    Max Vector3
}
\end{lstlisting}

    \item \texttt{OctreeStats} digunakan untuk menyimpan data statistik selama proses pembentukan pohon, seperti jumlah \textit{node} yang terbentuk dan jumlah \textit{node} yang dipangkas (\textit{pruned}) pada setiap tingkat kedalaman.
\begin{lstlisting}[language=Go, basicstyle=\ttfamily\small, breaklines=true]
type OctreeStats struct {
    CreatedNodesByDepth []int
    PrunedNodesByDepth  []int
    TotalVoxelLeaves    int
    mu                  sync.Mutex
}
\end{lstlisting}
\end{itemize}

Berikut adalah langkah-langkah algoritmanya:
\begin{enumerate}
    \item Program menghitung \textit{Bounding Box} utama dari seluruh \textit{vertex} objek untuk menentukan batas akar pohon.
    \item Fungsi rekursif \texttt{BuildOctree} dipanggil mulai dari kedalaman 0.
    \item Pada setiap pemanggilan, program mencatat statistik pembentukan \textit{node} pada kedalaman tersebut.
    \item Program membagi ruang saat ini menjadi 8 bagian menggunakan fungsi \texttt{Divide}.
    \item Untuk setiap bagian, program memeriksa apakah terdapat \textit{faces} (segitiga) objek yang berpotongan dengan bagian tersebut menggunakan algoritma \textit{Separating Axis Theorem} (SAT).
    \item Jika bagian tidak berpotongan dengan objek apa pun, bagian tersebut ditandai sebagai kosong (\textit{pruned}) dan rekursi berhenti untuk cabang tersebut.
    \item Jika bagian berpotongan dan belum mencapai kedalaman maksimal, program melakukan rekursi ke oktan tersebut.
    \item Jika sudah mencapai kedalaman maksimal (\textit{leaf}) dan terdapat potongan objek, bagian tersebut ditandai sebagai \textit{solid voxel}.
\end{enumerate}

Algoritma ini memiliki kompleksitas yang bergantung pada kedalaman pohon ($D$) dan jumlah \textit{faces} ($F$), secara teoritis sebesar $O(8^D \cdot F)$ untuk kasus terburuk, namun penggunaan pemangkasan (\textit{pruning}) secara signifikan mempercepat eksekusi pada area kosong.

\subsection{Algoritma Utama dalam Go}\label{sub:Algoritma Utama dalam Go} % (fold)

Implementasi fungsi pembangunan pohon secara paralel menggunakan \textit{concurrency} di Go:
\begin{lstlisting}[language=Go, basicstyle=\ttfamily\small, breaklines=true]
func BuildOctreeConcurrent(node *OctreeNode, vertices []Vector3, faces []Face, currentDepth, maxDepth int, stats *OctreeStats, wg *sync.WaitGroup) {
    defer wg.Done()
    // Update statistik dengan mutex lock
    stats.mu.Lock()
    stats.CreatedNodesByDepth[currentDepth]++
    stats.mu.Unlock()

    if currentDepth == maxDepth {
        node.IsLeaf = true
        if len(faces) > 0 {
            node.IsSolid = true
            // Mencatat leaf yang solid
        }
        return
    }

    childBoxes := Divide(node.Boundary)
    var childWg sync.WaitGroup
    for i := 0; i < 8; i++ {
        // Filter faces yang berpotongan (D&C)
        if len(intersectingFaces) == 0 {
            // Pruning node kosong
        } else {
            childWg.Add(1)
            go BuildOctreeConcurrent(childNode, vertices, intersectingFaces, currentDepth+1, maxDepth, stats, &childWg)
        }
    }
    childWg.Wait()
}
\end{lstlisting}

\subsection{Fungsi Pembantu}\label{sub:Fungsi Pembantu} % (fold)
Untuk mendukung berjalannya program utama, digunakan beberapa fungsi pembantu:
\begin{enumerate}
    \item \texttt{IsIntersectingSAT()}. Fungsi ini mengimplementasikan \textit{Separating Axis Theorem} untuk mendeteksi perpotongan antara segitiga (objek) dengan kotak (\textit{voxel}) secara akurat.
    \item \texttt{ParseOBJ()}. Fungsi ini bertugas membaca berkas masukan \texttt{.obj}, mengekstraksi data \textit{vertex} dan \textit{face}, serta menghitung ukuran \textit{bounding box} awal.
    \item \texttt{CollectVoxels()}. Fungsi rekursif yang menelusuri pohon \textit{octree} untuk mengumpulkan seluruh \textit{boundary} dari simpul daun yang berstatus \textit{solid}.
    \item \texttt{ExportToObj()}. Fungsi untuk menuliskan data koordinat \textit{voxel} yang terkumpul kembali ke dalam format berkas \texttt{.obj} untuk divisualisasikan.
    \item \texttt{IsInternalFace()}. Fungsi pembantu pada modul optimasi untuk mengecek apakah suatu sisi \textit{voxel} bersentuhan dengan \textit{voxel} lain, sehingga sisi tersebut tidak perlu digambar (optimasi visualisasi).
\end{enumerate}
